% MOVED OUT OF THE MAIN TEXT on 28 Oct for the ten-page limit (the author's decision;
% notes/22_page_cut_plan.md). This file is now input by appendix.tex, not by main.tex,
% which is why its heading is \section and its main-text references are numbers.
% Privacy-first section 7: what a practitioner can do, and why the usual metric will not help.
% Budget 1.25 pages including Table~\ref{tab:blind} and Figure~\ref{fig:blind}.
% Early stopping: outputs/nt_{ckpt,seeds}/nt250m_mlm_ecoli_s*, ep3/ep15/ep50, one basis 19 Oct.
% Blind metric: outputs/nt_ckpt/nt250m_mlm_human_s0/utility_at_leak.json and bpe_human/*, 19-20 Oct.
\section{What a practitioner can do, and what will not warn them}
\label{sec:blind}

Our measurements inform checkpoint selection in two ways: early checkpoints can have
lower measured recovery, and downstream task scores need not track that change. The comparisons below
are specific to the model, checkpoint and copy tier; they do not establish early stopping as a general defense.

\paragraph{Early checkpoints can resemble controls on the tested diagnostics.}
At round~3 of NT~250M on \emph{E.~coli}, seeds~1 and~2 have zero prefix and re-masking
recovery at both one and sixteen copies (50 records per tier). Single-copy target surprise
is $2.168$ and $2.154$\,\bits{}, against $2.112$ and $2.110$ for the respective controls;
LOO top-1 accuracy is $0.00$ after rounding, against $0.01$ for controls. At sixteen copies,
target surprise is $2.128$ and $2.115$\,\bits{} and LOO top-1 is $0.01$ in both seeds.
These are token-level conditional probes, not tests that a record is absent. The near-control
diagnostics and zero recovery describe these checkpoints under the tested protocols.

For records planted sixteen times, a further increase of three hundredths of a bit
in held-out loss yields 8\%, 8\% and 14\% exact recovery across three seeds. At a larger increase of +0.44\,\bits, the recovery rate reaches 66--76\%. Earlier stopping reduces recovery in this arm, with a sharp change between saved checkpoints. Our own from-scratch backbone tolerates roughly ten times that drift before
showing exact recovery under the tested protocol, so a study of that backbone alone reads more reassuringly than
released weights warrant. That is why this paper is built on released weights.

\paragraph{The downstream selection metric changes in both directions.}
Checkpoint selection can instead use downstream task scores.
Table~\ref{tab:blind} puts that score beside the leak at selected checkpoints of the same runs.

\begin{table}[tbp]
\centering
\caption{Downstream task score and leak measured on the same weights at selected checkpoints.
GUE promoter classification, linear probe on frozen mean-pooled states, Matthews correlation on the
full held-out split. \emph{id} is single-copy identification in the record's own host window, as in
Table~\ref{tab:ident}; an earlier version of this table quoted the average over the two host
conditions, which is a different quantity. Identification is over three seeds. DNABERT-2 MCC
entries give the three-seed minimum--maximum; NT MCC is from seed~0.
The DNABERT-2 best-checkpoint label uses mean linear-probe MCC across seeds.}
\label{tab:blind}
\begin{tabular}{llcccc}
\toprule
model & checkpoint & held-out \bits & MCC & 16-copy verbatim & 1-copy id \\
\midrule
NT 250M, chr22 & round 15 (best MCC) & 1.411 & \textbf{0.802} & \textbf{6\%}  & \textbf{0.04--0.08} \\
               & round 30            & 1.501 & 0.798          & 10\%          & 0.30--0.34 \\
               & round 50            & 1.632 & 0.780          & 18\%          & 0.60--0.66 \\
\midrule
DNABERT-2, chr22 & \shortstack[l]{round 15\\(best mean MCC)} & 1.264 & 0.798--0.801 & 0--2\%   & 0.04 \\
                 & round 120           & 1.600 & 0.784--0.792 & 42--46\% & 0.74--0.80 \\
                 & final               & 1.729 & 0.782--0.787 & 40--44\% & 0.92--0.96 \\
\bottomrule
\end{tabular}
\end{table}

On DNABERT-2 the score declines across the run by 1.1--1.6 points of MCC in total, monotonically in
two seeds and with one reversal at round~50 in the third, while single-copy identification climbs
from 4\% at round~15 to 92--96\% at the final checkpoint in its own window
(Figure~\ref{fig:blind}); matched own-window controls at round~15 are 0--2\%.
Selecting by mean linear-probe MCC across the three seeds chooses round~15, where measured
own-window membership rates are lower than at later checkpoints. Individual-seed selection
chooses round~15 for seeds~0--1 and round~50 for seed~2. This linear-probe score is not
a privacy signal: across the entire interval in which the model learns to disclose
membership, a practitioner watching it sees about a point of MCC of drift, which is what an
ordinary run looks like.

On NT fine-tuned on chromosome~22 the situation is worse, because the metric does not select a checkpoint with zero measured leakage. The highest task score belongs to the earliest checkpoint saved, and that one
already reproduces 6\% of repeated records in seed~0; single-copy own-window
rank-1 rates are 4--8\% across the three seeds. Choosing the best
available model by downstream performance means choosing a model that is already leaking.

\begin{figure}[!b]
\centering
\includegraphics[width=\linewidth]{F7_blind.pdf}
\caption{Downstream score and leak, measured on the same weights at the same checkpoints.
\textbf{Left:} DNABERT-2 on chromosome~22: the score falls by 1.1--1.6 points of MCC while
single-copy own-window rank-1 rates rise from 4\% at round~15 to 92--96\% at the final
checkpoint; round-15 matched controls are 0--2\%. Solid and dashed curves show the three-seed
means of MCC and rank-1, respectively; shading spans minimum--maximum MCC. The marked best
checkpoint maximizes mean linear-probe MCC.
\textbf{Right:} Nucleotide Transformer 250M on the same corpus, seed~0. The best downstream
score is at the earliest checkpoint shown, where single-copy own-window rank-1 is 8\%. Linear
probe on frozen mean-pooled states, GUE promoter classification, full held-out split.}
\label{fig:blind}
\end{figure}



\paragraph{A larger background corpus dilutes the leak, and does not prevent it.}
The obvious question about every number above is whether a real cohort, which is larger than ours,
is safer by being larger. Our experiment varies the background corpus rather than the participant cohort: the same backbone, cut
size, objective, canaries and schedule, with the fine-tuning corpus at 1{,}000, 2{,}500, 5{,}000,
10{,}000 and 20{,}000 windows of chromosome~22, two seeds each. The number of participants, how
many windows each contributes and how closely they are related are fixed by construction and were
not manipulated, so this experiment measures window-count scaling and bears on cohort size only to the
extent that the two move together. Table~\ref{tab:scale} reports results at matched held-out loss. This is necessary because, at a
fixed epoch count, a larger corpus also entails more optimization updates.

Dilution is real and it is not a defense. Twenty times the corpus divides single-copy verbatim
recovery by about four at $2.05$\,\bits{} (97\% to 23\%) and takes identification from $1.00$ to
$0.64$. At $2.20$\,\bits{}, all three corpus sizes that reach that loss have rank-1 rates
of $0.91$ or better; at $2.40$, every size that reaches it is at $1.00$. A larger corpus shifts the onset of the same leakage along the loss axis. This delay is useful and
uses the same measure as \S\ref{sec:leak}, but it does not provide immunity. At each corpus size's own held-out optimum, meanwhile,
single-copy verbatim recovery is 0\% at every size and sixteen-copy recovery is 78--100\% at every
size: the stopping rule does the work, not the size of the corpus.

\begin{table}[tbp]
\centering
\caption{Identification and verbatim recovery for one planted record as the background corpus grows,
read at matched held-out loss, mean of two seeds, records planted once. The manipulated quantity is the number of
chromosome~22 windows, not the number of participants. A dash means that corpus never reaches that
loss in 30 epochs. \emph{Best loss} is the corpus's own optimum, which improves with size as it
should.}
\label{tab:scale}
\begin{tabular}{lccccc}
\toprule
windows & best loss & id @2.05 & id @2.20 & verbatim @2.05 & verbatim @2.20 \\
\midrule
1{,}000  & 1.954 & 1.00          & ---           & 97\%          & --- \\
2{,}500  & 1.910 & 1.00          & ---           & 96\%          & --- \\
5{,}000  & 1.883 & 0.95          & 1.00          & 56\%          & 94\% \\
10{,}000 & 1.789 & 0.74          & 1.00          & 50\%          & 82\% \\
20{,}000 & 1.708 & \textbf{0.64} & \textbf{0.91} & \textbf{23\%} & \textbf{37\%} \\
\bottomrule
\end{tabular}
\end{table}

\paragraph{Recommendation.}
These findings inform four practical recommendations, presented as guidance distinct from the
experimental results themselves.

\emph{Stop where held-out loss stops improving, and compute it.} In this trajectory, held-out loss worsens before the observed rise in exact
recovery, and it is inexpensive to monitor during fine-tuning. This stopping point does not define a sharp boundary between privacy and leakage. At that checkpoint single-copy identification is $0.04$ and exact recovery is zero,
but the decoder's output is already closer to a planted single-copy record than to a never-trained
one in the seed-averaged descriptive comparison ($31.7$ against $33.2$ nucleotides
wrong out of 48); the within-cell gaps are small relative to their standard errors
(\S\ref{sec:ident}). The earlier checkpoint has lower measured leakage under the tested protocols;
this does not establish that its records are unreachable by another attack. It is a measured
risk reduction in these runs, not a safety property of a release: at that same checkpoint a
record planted sixteen times is already ranked first.

\emph{Do not treat downstream score as a safety signal.} Between each model's own held-out optimum
and the end of its run, an interval over which extraction goes from 0--6\% to 36.7--100\%, the available frozen-probe means
on six tasks across two organisms show maximum absolute changes of $0.024$ MCC on the four DNABERT
arms and $0.083$ on the other two, while the fine-tuned head moves \emph{upward}
on four of the six tasks for DNABERT-2, telling a practitioner who watches it to keep training.
Separately, the frozen promoter probe on chromosome~22 declines by 1.1--1.6 MCC points
as membership rank-1 rates rise. Across these evaluations, downstream utility does not provide
a consistent warning of the observed leakage.

\emph{Evaluate masking under the intended model and release conditions.} On our own
backbone, trained on the same corpus for the same budget, replacing next-token training with masked
training takes verbatim extraction from 100\% to \emph{zero} in every cell we ran (four corpora,
two mask rates, three cut sizes), and it \emph{improves} held-out loss, from 2.3--4.2\,\bits{} to
1.66--1.98. This is a 100-percentage-point reduction in exact extraction within the
controlled objective comparison, for the tested backbone and decoding protocols.
The masked pretrained models in Table~\ref{tab:cost} also exhibit leakage, which limits
transfer of this observation. That comparison changes architecture, scale and training
history together; it does not isolate pretraining as the cause or establish masking as
a privacy defense. Both leakage endpoints require evaluation on the intended released model.

\emph{Evaluate tokenization within the intended backbone and objective.} Under the drawn
placement used in the controlled sweep, a model whose canary occupies 24 or 48 content slots
returned no record whole at any loss we drove it to, including losses far past usefulness. That
sweep measured exact extraction as its endpoint; whether more slots also reduce membership
disclosure was not tested there. The placement condition is also essential to this interpretation: at a \emph{fixed} offset the character tokenizer reaches 100\% at $2.20$\,\bits{}
(\S\ref{sec:basis}, Table~\ref{tab:capacity}). The effect therefore depends on placement,
and the differing target surprise does not support equal-learning comparisons. HyenaDNA
and the failed overlapping-family prediction further limit a general slot-count explanation
(\S\ref{sec:slots}). The magnitude remains substantial: exact recovery rises from 0\% with 2-mers
to 92\% with 6-mers at matched loss $2.40$\,\bits{} in the sixteen-copy sweep. Fine-grained
tokenizers carry their own costs in sequence length and downstream accuracy, which depend on the task and were not evaluated comprehensively here. Slot count is therefore a
design choice whose costs merit measurement; we do not recommend it without qualification.
